\documentclass[11pt]{article}
\usepackage[a4paper,margin=27mm]{geometry}
\usepackage[T1]{fontenc}
\usepackage{amsmath,amssymb,amsthm,booktabs,array}
\usepackage[hidelinks]{hyperref}

\newcommand{\F}{\mathbb F}
\newcommand{\Q}{\mathbb Q}
\newcommand{\Z}{\mathbb Z}
\newcommand{\lfd}{\operatorname{LFD}}
\newcommand{\NF}{\operatorname{NF}}
\newcommand{\wt}{\operatorname{wt}}
\newcommand{\Res}{\operatorname{Res}}
\newcommand{\code}[1]{\texttt{#1}}

\newtheorem{theorem}{Theorem}
\newtheorem{proposition}[theorem]{Proposition}
\newtheorem{lemma}[theorem]{Lemma}
\newtheorem{corollary}[theorem]{Corollary}
\theoremstyle{definition}
\newtheorem{definition}[theorem]{Definition}
\theoremstyle{remark}
\newtheorem{remark}[theorem]{Remark}

\title{The carry lower bound for every number of summands}
\date{Working note, 2026-11-20}

\begin{document}
\sloppy
\maketitle

\begin{abstract}
For the direct $m$-summand two-power-subgroup chain presentation
$G_{m,k}\cup\{\NF(S_{m+1}(x_1,\ldots,x_m,t))\}$ of the
Petit--Kosters--Messeng summation system on $y^2=x^3+ax+b$, we prove
\[
 \lfd\ \ge\ m\cdot\min(k,m-1)+1 ,
\]
for every $m\ge3$ and every depth $k\ge1$, over every field with $2\ne0$,
under explicit nondegeneracy hypotheses. For $k\ge m$ the only hypothesis
is $c_m(a,b,t)\ne0$, where $c_m$ is the coefficient of
$x_1^{2^{m-1}-1}\cdots x_m^{2^{m-1}-1}$ in $S_{m+1}(x_1,\ldots,x_m,t)$; it
generalizes $c_3=-8(5t^3+3at+6b)$ of the three-summand theorem. For
$k\le m-1$ the bound is the exact value $mk+1$ when the normal form has its
full degree $mk$ and $S_{m+1}(\cdot,t)$ is nonzero at two or more points
of the grid $\mu_{2^k}^m$ (for instance on every unit ideal). We show that $c_m$ has the form
$\kappa_m t^m+(\text{terms with }a\text{ or }b)$ with
$\kappa_m=[s_1^2\cdots s_m^2]\prod_{\varepsilon\in\{\pm1\}^m}(1+\varepsilon\cdot
s)\ne0$, so for every curve at most $m$ targets are excluded (over $\Q$, and
over $\F_p$ whenever $p\nmid\kappa_m$). We list
$c_2,\ldots,c_5$ explicitly, validate the statements against every
measured value, and formalize the abstract core for every $m$ in Lean.
\end{abstract}

\paragraph{Status.} Every statement below has a complete written proof.
The abstract theorems (Theorems~\ref{thm:onset} and~\ref{thm:full} for an
arbitrary generator, and the three-summand instance) are formalized in
\code{formal/ICFallDegrees/CarryAllM.lean}, which builds sorry-free against the pinned
Lean and Mathlib (Section~\ref{sec:lean}). The facts about $S_{m+1}$ in
Section~\ref{sec:summation} (the shape of $c_m$, the identity on the
cuspidal cubic and $\kappa_m\ne0$) are written proofs, checked numerically
for $m\le5$ (Section~\ref{sec:numerics}); they are not
formalized.

\section{Setting}
Let $K$ be a field with $2\ne0$, $m,k\ge1$. The chain ring has variables
$u_{i,j}$, $1\le i\le m$, $0\le j<k$, and the chain equations
\[
 G_{m,k}=\{u_{i,j}^2-u_{i,j+1}\ (j+1<k),\quad u_{i,k-1}^2-1\}.
\]
A polynomial is \emph{reduced} if all its monomials are squarefree. By the
kernel-checked chain theorem of the paper (Theorem \emph{Chain subsystem}, label \code{thm:reduced},
\code{ChainQuotient.reduced\_chain\_ideal\_trivial}) a reduced polynomial in
$\langle G_{m,k}\rangle$ is zero, so every polynomial $f$ has a unique reduced
$\NF(f)\equiv f$ modulo the chains. For an exponent $e\in\Z_{\ge0}$ write
$\beta_j(e)\in\{0,1\}$ for its binary digits and
$\wt_k(e)=\sum_{j<k}\beta_j(e)$.

For $S\in K[x_1,\ldots,x_m]$ we write $\NF(S)$ for the normal form of
$S(u_{1,0},\ldots,u_{m,0})$. The direct presentation of $m$ summands with
target abscissa $t$ is $F_{m,k}=G_{m,k}\cup\{g\}$, $g=\NF(S_{m+1}(x_1,\ldots,x_m,t))$.
The last fall degree $\lfd(F)$ is that of~\cite{hkyy} as in the paper:
the least $d$ such that every $h\in\langle F\rangle$ lies in the degree-$\max(d,\deg h)$ closure
$V_{F,\max(d,\deg h)}$.

\begin{lemma}[Normal form of a monomial]\label{lem:nf}
For every exponent vector $e$ and $k\ge1$,
\[
 \NF\Bigl(\prod_i x_i^{e_i}\Bigr)=\prod_{i=1}^m\prod_{j<k}u_{i,j}^{\beta_j(e_i)} ,
\]
a squarefree monomial of degree $\sum_i\wt_k(e_i)$. On exponents with all
$e_i<2^k$ the map $e\mapsto\NF(x^e)$ is injective.
\end{lemma}
\begin{proof}
Modulo the chains $u_{i,0}^{2^j}\equiv u_{i,j}$ for $j<k$ (induction on $j$
with $u_{i,j}^2\equiv u_{i,j+1}$) and $u_{i,0}^{2^k}\equiv u_{i,k-1}^2\equiv1$.
Hence $u_{i,0}^{e}\equiv u_{i,0}^{e\bmod 2^k}=\prod_{j<k}(u_{i,0}^{2^j})^{\beta_j(e)}
\equiv\prod_{j<k}u_{i,j}^{\beta_j(e)}$, since the digits $j<k$ of $e$ and of
$e\bmod2^k$ agree. The right side is squarefree, hence the normal form. A
number below $2^k$ is determined by its digits below $k$.
\end{proof}

\begin{lemma}[Carry lemma; kernel-checked]\label{lem:carry}
Let $g$ be reduced of degree $d$ and $h$ reduced of degree at most $d$ with
$h\equiv A g$ modulo the chains for some polynomial $A$, and $h\ne cg$ for
every $c\in K$. Then $\lfd(G_{m,k}\cup\{g\})\ge d+1$.
\end{lemma}
This is the first statement of the paper's carry theorem (\code{thm:carry})
(\code{CarryLower.chain\_carry\_lower}). Proof in one line: $V_d$ lies in
$(\langle G_{m,k}\rangle\cap M_{\le d})+Kg$, because a row of degree $<d$
in that space is already in $\langle G_{m,k}\rangle$; the ideal element $h$
of degree $\le d$ would then satisfy $h-cg\in\langle G_{m,k}\rangle$ with
$h-cg$ reduced, so $h=cg$.

\begin{lemma}[Digit weights]\label{lem:weights}
Let $n\le k$. Then $\wt_k(2^n-1)=n$, and every $x<2^n$ with $x\ne2^n-1$ has
$\wt_k(x)\le n-1$. For $m\ge3$ every $x\le2^{m-1}+1$ satisfies
$x<2^m-1$, hence $\wt_k(x)\le m-1$ when $k\ge m$.
\end{lemma}
\begin{proof}
Digits at or above $n$ of $x<2^n$ vanish, so $\wt_k(x)=\wt_n(x)$. If all
$n$ low digits of $x$ are $1$ then $x=2^n-1$. For $m\ge3$,
$2^m-1-(2^{m-1}+1)=2^{m-1}-2>0$.
\end{proof}

\section{The two bounds for an arbitrary generator}

\begin{theorem}[Onset bound; kernel-checked as \code{CarryAllM.lfd\_ge\_onset}]\label{thm:onset}
Let $m\ge3$, $k\ge m$, and let $S\in K[x_1,\ldots,x_m]$ have degree at most
$2^{m-1}$ in each variable, with a nonzero coefficient at
$\mathbf x^{\tau}=x_1^{2^{m-1}-1}\cdots x_m^{2^{m-1}-1}$. Then
$g=\NF(S)$ has degree exactly $m(m-1)$, and
\[
 \lfd\bigl(G_{m,k}\cup\{g\}\bigr)\ \ge\ m(m-1)+1 .
\]
\end{theorem}
\begin{proof}
Take $h=\NF(x_1S)\equiv u_{1,0}\,g$ and apply Lemma~\ref{lem:carry} with
$A=u_{1,0}$ and $d=m(m-1)$. Three facts are needed.

\emph{$\deg g=m(m-1)$.} Every exponent of $S$ has $e_i\le2^{m-1}$, so by
Lemmas~\ref{lem:nf} and~\ref{lem:weights} every monomial of $g$ has
degree $\sum_i\wt_k(e_i)\le m(m-1)$. All exponents are below $2^{m-1}+1<2^k$,
so the monomials $\NF(x^e)$ are distinct and the coefficient of
$\NF(\mathbf x^\tau)$ in $g$ equals $[\mathbf x^\tau]S\ne0$; that monomial has
degree $m\,\wt_k(2^{m-1}-1)=m(m-1)$ because $m-1\le k$.

\emph{$\deg h\le m(m-1)$.} Exponents of $x_1S$ satisfy $e_1\le2^{m-1}+1$ and
$e_i\le2^{m-1}$ otherwise, so each block has weight at most $m-1$ by
Lemma~\ref{lem:weights} (this is where $m\ge3$ and $k\ge m$ enter).

\emph{$h\ne cg$.} All exponents of $S$ and $x_1S$ are below $2^m\le2^k$, so
$\NF$ is injective on their supports (Lemma~\ref{lem:nf}) and $h=cg$ would give
$x_1S=cS$ as polynomials, i.e.\ $(x_1-c)S=0$, so $S=0$, contradicting
$[\mathbf x^\tau]S\ne0$.
\end{proof}

\begin{theorem}[Full-degree bound; kernel-checked as \code{CarryAllM.lfd\_ge\_full}]\label{thm:full}
Let $m,k\ge1$, $S\in K[x_1,\ldots,x_m]$ arbitrary and $g=\NF(S)$. If
$\deg g=mk$ and $\NF(x_iS)$ is not a scalar multiple of $g$ for some $i$,
then $\lfd(G_{m,k}\cup\{g\})\ge mk+1$. If moreover $\deg g\ge2$, then
$\lfd=mk+1$ exactly.
\end{theorem}
\begin{proof}
$h=\NF(x_iS)\equiv u_{i,0}g$ is reduced, so $\deg h\le mk$; apply
Lemma~\ref{lem:carry} with $d=mk$. The upper bound is the kernel-checked
$\lfd\le\max(mk+1,\delta)$ of the paper (\code{cor:allm}).
\end{proof}

\begin{proposition}[When the second hypothesis holds]\label{prop:grid}
Suppose $\deg g=mk$, let $\bar K$ be an algebraic closure of $K$ and
$H=\mu_{2^k}(\bar K)$, a group of order $2^k$. If $S$ is nonzero at two
distinct points of $H^m$ (equivalently, the system has at most $2^{mk}-2$
solutions on $H^m$; for instance any unit ideal), then $\NF(x_iS)\notin Kg$
for some $i$, and $\lfd=mk+1$.
\end{proposition}
\begin{proof}
Normal forms commute with the extension $K\subset\bar K$ (they are
computed by Lemma~\ref{lem:nf}), so it suffices to work over $\bar K$.
Under $u_{i,j}\mapsto x_i^{2^j}$, reduced polynomials correspond bijectively to
polynomials of degree $<2^k$ in each $x_i$ (binary digits), and these are
determined by their values on $H^m$ (grid vanishing; $|H|=2^k$ because
$2\ne0$). A polynomial and its normal form agree on $H^m$, since the chain
equations vanish there. If $\NF(x_iS)=c_ig$ for every $i$, then
$(x_i-c_i)S$ vanishes on $H^m$ for every $i$. Let $P\ne Q$ be points of
$H^m$ with $S(P)S(Q)\ne0$ and choose $i$ with $P_i\ne Q_i$; then
$c_i=P_i$ and $c_i=Q_i$, a contradiction.
\end{proof}

\section{The summation polynomial}\label{sec:summation}
Define $S_3$ as in the paper and, for $n\ge3$,
\[
 S_{n+1}(x_1,\ldots,x_n,x_{n+1})=\Res_X\bigl(S_n(x_1,\ldots,x_{n-1},X),\,S_3(x_n,x_{n+1},X)\bigr),
\]
the Sylvester resultant with formal degrees $2^{n-2}$ and $2$. Then
$S_{n+1}\in\Z[a,b][x_1,\ldots,x_{n+1}]$ has degree at most $2^{n-1}$ in each
variable (by induction: the resultant is homogeneous of degree $2$ in the coefficients of
its first argument and of degree $2^{n-2}$ in those of its second, and $S_3$ is quadratic in each argument). Other orders of the resultants
differ from it at most by a nonzero constant factor (both are irreducible with the same zero set~\cite{semaev}), which does not
affect any hypothesis below; the numerics of Section~\ref{sec:numerics} use this recursion, which
agrees with the paper's $S_4$ exactly.

\begin{definition}
For $m\ge2$ let $\mathbf x^\tau=\prod_{i\le m}x_i^{2^{m-1}-1}$ and
\[
 c_m(a,b,t)=[\mathbf x^\tau]\,S_{m+1}(x_1,\ldots,x_m,t)\in\Z[a,b,t],
\qquad
 \gamma_{m,k}(a,b,t)=\sum_{\substack{0\le e_i\le2^{m-1}\\ e_i\equiv-1\ (2^k)}}[\mathbf x^e]\,S_{m+1}(x,t).
\]
\end{definition}
Since $2^{m-1}-1$ is the only exponent $\le2^{m-1}$ of weight $m-1$, for
$k\ge m-1$ we have $\deg\NF(S_{m+1})=m(m-1)$ if and only if $c_m\ne0$, and
$\gamma_{m,m-1}=c_m$. For $k\le m-1$ the coefficient of the full monomial
$\prod u_{i,j}$ in $g$ is $\gamma_{m,k}$ (Lemma~\ref{lem:nf}), so
$\deg g=mk$ if and only if $\gamma_{m,k}\ne0$.

\paragraph{Weights.} Give $x_i,t$ weight $1$, $a$ weight $2$ and $b$ weight $3$.
The scaling $(x,y)\mapsto(\lambda^2x,\lambda^3y)$, $(a,b)\mapsto(\lambda^4a,\lambda^6b)$
shows that $S_3$ is weighted homogeneous of weight $4$; a resultant of
weighted homogeneous polynomials of weights $w_A,w_B$ and $X$-degrees
$n_A,n_B$ ($X$ of weight $1$) is weighted homogeneous of weight
$n_B(w_A-n_A)+n_Aw_B$, so $S_{m+1}$ has weight $m2^{m-1}$ by induction.
Hence $c_m$ is weighted homogeneous of weight $m$ in $(t,a,b)$, and its
coefficient of $t^m$ is its value at $a=b=0$, $t=1$. Explicitly (computed
with exact interpolation over $\F_{2^{31}-1}$ at random $(a,b,t)$, solved on a weighted-monomial basis and
checked at two further points; Section~\ref{sec:numerics}):
\begin{align*}
 c_2&=-2\,(t^2+a),\\
 c_3&=-8\,(5t^3+3at+6b),\\
 c_4&=-16\,(47t^4+84at^2-12bt-3a^2),\\
 c_5&=-128\,(1487t^5-365at^3+1545bt^2+318a^2t+57ab).
\end{align*}
$c_2$ is checked by hand and $c_3$ is the kernel-checked
\code{M3CarryLower.coeff\_s4\_333}.

\begin{proposition}[The cuspidal cubic]\label{prop:cusp}
Let $N_m(s_0;s_1,\ldots,s_m)=\prod_{\varepsilon\in\{\pm1\}^m}(s_0+\varepsilon_1s_1+\cdots+\varepsilon_ms_m)\in\Z[s]$.
At $a=b=0$, for every $m\ge2$,
\[
 \Bigl(\prod_{i=0}^m s_i^{2^m}\Bigr)\,S_{m+1}\bigl(s_1^{-2},\ldots,s_m^{-2},s_0^{-2}\bigr)=N_m(s_0;s_1,\ldots,s_m).
\]
Consequently $c_m(0,0,t)=\kappa_m t^m$ with
$\kappa_m=[s_1^2\cdots s_m^2]\,N_m(1;s_1,\ldots,s_m)$.
\end{proposition}
\begin{proof}
For $m=2$ both sides equal
$s_0^4+s_1^4+s_2^4-2s_0^2s_1^2-2s_0^2s_2^2-2s_1^2s_2^2$ (direct expansion;
$N_2=((s_0+s_1)^2-s_2^2)((s_0-s_1)^2-s_2^2)$).
Assume the identity for $m$. Fix representatives $\varepsilon_c$ of the
$2^{m-1}$ classes $c$ of $\{\pm1\}^m/\pm$ and write $\ell_c=\varepsilon_c\cdot s$.
Pairing $\varepsilon$ with $-\varepsilon$, $N_m(\sigma;s)=\prod_c(\sigma^2-\ell_c^2)$, so
with $x_i=s_i^{-2}$ the hypothesis reads, as an identity of polynomials in $X$,
\begin{align*}
 A(X):=S_{m+1}(x_1,\ldots,x_m,X)&=\prod_{i\le m}x_i^{2^{m-1}}\prod_c(1-\ell_c^2X)
 =L\prod_c\bigl(X-\ell_c^{-2}\bigr),\\
 L&=\prod_{i\le m}x_i^{2^{m-1}}\prod_c(-\ell_c^2).
\end{align*}
With $B(X)=S_3(x_{m+1},x_0,X)$ and the case $m=2$ (and the symmetry of $S_3$),
$B(\ell_c^{-2})=N_2(\ell_c;s_{m+1},s_0)\,s_{m+1}^{-4}s_0^{-4}\ell_c^{-4}$. The resultant
formula $\Res_X(A,B)=L^2\prod_cB(\ell_c^{-2})$ holds over $\Q(s)$ because
$L\ne0$. The $2^{m+1}$ linear forms $\ell_c+\delta s_{m+1}+\delta's_0$
($c$, $\delta,\delta'\in\{\pm1\}$) are $\delta'(s_0+\varepsilon''\cdot(s_1,\ldots,s_{m+1}))$
with $\varepsilon''$ running once through $\{\pm1\}^{m+1}$, and each sign
$\delta'$ occurs $2^m$ times, so $\prod_cN_2(\ell_c;s_{m+1},s_0)=N_{m+1}(s_0;s_1,\ldots,s_{m+1})$.
Multiplying out,
$\Res_X(A,B)=N_{m+1}\big/\bigl(s_0^{2^{m+1}}\prod_{i=1}^{m+1}s_i^{2^{m+1}}\bigr)$,
the identity for $m+1$.

For the coefficient: the monomial $\prod_{i=0}^ms_i^{2d_i}$ of $N_m$
($\sum d_i=2^{m-1}$; $N_m$ is even in each $s_i$ and homogeneous of degree $2^m$)
corresponds to $\prod_{i\ge1}x_i^{2^{m-1}-d_i}\,t^{2^{m-1}-d_0}$. The exponent
$\tau$ is $d_1=\cdots=d_m=1$, which forces $d_0=2^{m-1}-m$ and the power
$t^m$; dehomogenizing at $s_0=1$ gives $\kappa_m$.
\end{proof}

\begin{proposition}[$\kappa_m\ne0$]\label{prop:kappa}
For $m\ge1$,
\[
 \kappa_m=2^{-m}\sum_{\pi}\ (-2^m)^{|\pi|}\prod_{B\in\pi}(2|B|-1)! ,
\]
the sum over the set partitions $\pi$ of $\{1,\ldots,m\}$. For $m\ge2$,
$v_2(\kappa_m)=v_2((2m-1)!)=2m-1-s_2(2m-1)$, where $s_2$ is the binary
digit sum; in particular $\kappa_m\ne0$.
The first values are $\kappa_2=-2$, $\kappa_3=-40=-2^3\cdot5$,
$\kappa_4=-752=-2^4\cdot47$, $\kappa_5=-190336=-2^7\cdot1487$,
$\kappa_6=255410944=2^8\cdot997699$, $\kappa_7=-2^{10}\cdot13^2\cdot9470731$,
$\kappa_8=2^{11}\cdot4969\cdot3677156909$.
\end{proposition}
\begin{proof}
Put $y_i=s_i^2$ and $P=N_m(1;s)=\prod_\varepsilon(1+\varepsilon\cdot s)\in1+s\Q[[s]]$.
Averaging over $\varepsilon$ kills the odd powers:
$\log P=\sum_{n\ge1}\frac{(-1)^{n+1}}n\sum_\varepsilon(\varepsilon\cdot s)^n
=-2^m\sum_{j\ge1}\frac1{2j}\sum_{|\alpha|=j}\frac{(2j)!}{\prod(2\alpha_i)!}y^\alpha$.
The coefficient of the squarefree monomial $\prod_{i\in B}y_i$ in $\log P$ is
$-2^m(2|B|-1)!/2^{|B|}$, and the coefficient of $y_1\cdots y_m$ in
$P=\exp(\log P)$ is the sum over set partitions of the products of these block
coefficients, which is the formula.

In $2^m\kappa_m=\sum_\pi(-2^m)^{|\pi|}\prod(2|B|-1)!$ the partition into $q$
blocks of sizes $j_1,\ldots,j_q$ contributes a term of $2$-adic valuation
$mq+\sum_l v_2((2j_l-1)!)=mq+2m-q-\sum_ls_2(2j_l-1)$ (Legendre). For the
single block ($q=1$) this is $3m-1-s_2(2m-1)$. The difference for $q\ge2$ is
\[
 (m-1)(q-1)+s_2(2m-1)-\sum_l s_2(2j_l-1)\ \ge\ (m-1)(q-2)+s_2(2m-1)-1 ,
\]
using $s_2(2j-1)\le j$ ($2j-1<2^j$) and $\sum j_l=m$. For $q\ge3$ this is at
least $m-1\ge1$; for $q=2$ it is at least $s_2(2m-1)-1\ge1$, because $2m-1\ge3$
is odd and not a power of two. So the single-block term has strictly
smallest valuation and determines $v_2(2^m\kappa_m)$.
\end{proof}

\begin{corollary}[The hypotheses are generic]\label{cor:generic}
Let $m\ge3$ and $p$ the characteristic of $K$ ($p=0$ or $p\nmid2\kappa_m$).
\begin{enumerate}
\item For every $a,b\in K$, $c_m(a,b,t)$ is a polynomial of degree exactly $m$ in
$t$ with leading coefficient $\kappa_m$; so at most $m$ targets $t$ violate
$c_m\ne0$.
\item For every $1\le k\le m-1$, $\gamma_{m,k}\not\equiv0$: at $a=b=0$ the coefficient of
$t^m$ in $\gamma_{m,k}$ is $\kappa_m$.
\end{enumerate}
\end{corollary}
\begin{proof}
(1) By weighted homogeneity of weight $m$, $t^m$ is the only monomial of
$c_m$ free of $a,b$, and its coefficient is $c_m(0,0,1)=\kappa_m$.
(2) At $a=b=0$, $S_{m+1}$ is homogeneous of degree $m2^{m-1}$ in $(x,t)$,
so $[\mathbf x^e]S_{m+1}(x,t)$ is a multiple of $t^{m2^{m-1}-|e|}$. Exponents
with $e_i\equiv-1\pmod{2^k}$ and $e_i\le2^{m-1}$ satisfy $e_i\le2^{m-1}-1$,
with $|e|=m(2^{m-1}-1)$ only for $e=\tau$; so only $\tau$ contributes to
$t^m$, with $\kappa_m$ (Proposition~\ref{prop:cusp}).
\end{proof}

\section{The theorem for summation systems}
\begin{theorem}[Carry bound for every number of summands]\label{thm:main}
Let $K$ be a field with $2\ne0$, $m\ge3$, $a,b,t\in K$, $k\ge1$, and
$F_{m,k}=G_{m,k}\cup\{\NF(S_{m+1}(x_1,\ldots,x_m,t))\}$.
\begin{enumerate}
\item (Onset, $k\ge m$.) If $c_m(a,b,t)\ne0$, then $\lfd(F_{m,k})\ge m(m-1)+1$.
\item (Full degree, $k\le m-1$.) If $\gamma_{m,k}(a,b,t)\ne0$ and $S_{m+1}(\cdot,t)$ is
nonzero at two distinct points of $\mu_{2^k}(\bar K)^m$, then $\lfd(F_{m,k})=mk+1$.
\end{enumerate}
In both cases $\lfd(F_{m,k})\ge m\min(k,m-1)+1$, and this is
$\deg\NF(S_{m+1})+1$. At $k=m-1$ the two parts agree ($\gamma_{m,m-1}=c_m$).
By Corollary~\ref{cor:generic}, part~1 excludes at most $m$ targets for each
curve; in part~2 the condition $\gamma_{m,k}\ne0$ fails only on a proper
Zariski-closed subset of the $(a,b,t)$-space, and the grid condition holds
on every unit ideal (more generally whenever at most $2^{mk}-2$ grid points
are solutions).
\end{theorem}
\begin{proof}
Part~1 is Theorem~\ref{thm:onset} with $S=S_{m+1}(x,t)$, which has degree
$\le2^{m-1}$ in each $x_i$ and $[\mathbf x^\tau]S=c_m\ne0$. Part~2 is
Theorem~\ref{thm:full} with Proposition~\ref{prop:grid}; $\deg g=mk\ge3$.
\end{proof}

\begin{remark}[What the theorem does not give]
(i) It is a statement about the direct presentation; the split presentation
with auxiliary abscissae has smaller measured values (paper,
Table \code{tab:splitm}) and is not covered. (ii) For $k\ge m$ it is only a
lower bound; the upper bound is $mk+1$. The measured values exceed it at
$m=3$, $k\ge6$ ($8$), so the onset bound is not the growth law. (iii) On the
locus $c_m=0$ (for $k\ge m$) the theorem is silent; for $m=3$ the measured value
there is still $7=\deg g+2$ at $k=3,4$. (iv) For $k\ge m$ the multiplier
$u_{1,0}$ is essential: a larger multiplier can raise the degree of $h$
above $\deg g$.
\end{remark}

\section{Numerical validation}\label{sec:numerics}
\paragraph{Coefficients.} A C program in the scratch area (not committed;
the repository forbids new Python and this computation is a check, not a
result) evaluates $S_{m+1}$ by the recursion over $\F_p$ and extracts
$[\mathbf x^\tau]$ by tensor Lagrange interpolation on $2^{m-1}+1$ nodes per
variable. It reproduces $-8(5t^3+3at+6b)$ at five points; the four
(resp.\ five) weighted coefficients of $c_4$ (resp.\ $c_5$) were solved from
four (five) random points and confirmed at two more each. Its value of $c_m(0,0,1)$ equals the
$\kappa_m$ of the partition formula for $m=2,3,4,5$ ($-2,-40,-752,-190336$).
On the recorded
four-summand instances it gives $c_4\ne0$: U1 ($p=2013265921$) $c_4=431449230$
(the program and the closed form agree), U2 ($p=998244353$) $430513295$,
$P$-m4k3-p193 ($a=b=167$, $t=116$) $164$ (program and closed form agree), consistent
with the measured $\deg g=12$ (U1, U2 at $k=4,5$; the $p=193$ instance at $k=3$).

\paragraph{Measured values.} Every exact value or certified lower bound of
the paper and of the four-summand growth study, against the
prediction of Theorem~\ref{thm:main} (``pred.'').
\begin{center}\small
\begin{tabular}{ccp{2.1cm}p{6.0cm}p{2.6cm}}
\toprule
$m$ & $k$ & pred. & measured & agreement\\
\midrule
3 & 2 & $=7$ (part 2) & $7$ on $105/116$ and $48/48$ (the others have $c_3=0$, $\deg g=5$) & yes; silent where $c_3=0$\\
3 & 3,4,5 & $\ge7$ & $7$ ($234$ seeded, $12$ generic, U1--U3) & yes (tight)\\
3 & 3,4 on $c_3=0$ & silent & $7$ ($\deg g=5$) & ---\\
3 & 6 & $\ge7$ & $8$ (U1--U3, two proper instances) & yes (not tight)\\
3 & 7 & $\ge7$ & $\ge8$ & yes\\
4 & 2 & $=9$ & $9$ ($48/48$; six certified) & yes\\
4 & 3 & $=13$ & $13$ ($47/48$; the other has $\deg g=11$) & yes\\
4 & 4 & $\ge13$ & $13$ (six proper, U1, U2) & yes (tight)\\
4 & 5 & $\ge13$ & $\ge13$ (U1, U2) & yes\\
5 & 2,3,4 & $=11,16,21$ & $11,16,21$ & yes\\
6 & 2,3,4 & $=13,19,25$ & $13,19,25$ & yes\\
7 & 2,3 & $=15,22$ & $15,22$ ($47/48$ at $k=2$) & yes\\
8 & 2,3 & $=17,25$ & $17,25$ ($47/48$ at $k=2$) & yes\\
5--8 & $\ge m$ & $\ge m(m-1)+1$ & not computed ($>2^{24}$ grid points) & ---\\
\bottomrule
\end{tabular}
\end{center}
No measured value contradicts the theorem. The three exceptions of the
paper's Table \code{tab:msum} have $\deg g=mk-1$, i.e.\ $\gamma_{m,k}=0$, where the
theorem makes no claim. The paper's tables of $m$-summand values for
$k\le m-1$ were recorded as ``$\deg g=mk$ and the hypotheses on $h$ hold
computationally''; Theorem~\ref{thm:main}(2) replaces the second hypothesis by
the grid condition, which every unit instance satisfies.

\section{Lean status}\label{sec:lean}
\code{formal/ICFallDegrees/CarryAllM.lean} proves, for every $m$:
\code{lift\_congr} (Lemma~\ref{lem:nf} for every exponent),
\code{weight\_two\_pow\_sub\_one}, \code{weight\_le\_pred} (Lemma~\ref{lem:weights}),
\code{lfd\_ge\_onset} (Theorem~\ref{thm:onset}), \code{lfd\_ge\_full}
(Theorem~\ref{thm:full}, lower half) and \code{m3\_onset} (the three-summand
case recovered from the general theorem with \code{M3CarryLower.coeff\_s4\_333}).
Not formalized: Proposition~\ref{prop:grid}, Section~\ref{sec:summation}
(the definition of $S_{m+1}$ for general $m$, the multidegree bound, the
cusp identity and $\kappa_m\ne0$), and therefore the specialization of
\code{lfd\_ge\_onset} to $S_{m+1}$ for $m\ge4$.

The file was built with Lean 4.19.0 and Mathlib \code{v4.19.0} (the
repository pins) by \code{lake build ICFallDegrees.CarryAllM}; it contains no
\code{sorry}, \code{axiom} or \code{native\_decide}, and \code{\#print axioms} reports only
\code{propext}, \code{Classical.choice} and \code{Quot.sound} for every
declaration listed above. The module is imported by the root file and its
main declarations are added to \code{Audit.lean}; the full repository build
(all modules) was not rerun locally.

\begin{thebibliography}{9}
\bibitem{hkyy} M.-D.~Huang, M.~Kosters, Y.~Yang, S.~L.~Yeo, On the last fall degree of zero-dimensional Weil descent systems, J.~Symbolic Comput.\ 87 (2018).
\bibitem{semaev} I.~Semaev, Summation polynomials and the discrete logarithm problem on elliptic curves, IACR ePrint 2004/031.
\end{thebibliography}
\end{document}
